\begin{abstract}
Coding agents forget everything between sessions, and the memory systems built for them mostly select context by similarity top-$k$, ignore what items cost in tokens, and have no model of whether a stored fact is still true. We cast agent memory as online budgeted selection, show that the surveyed systems are special cases, and derive where each is worst-case suboptimal. We then present \kanerva, a local memory layer with no LLM in the loop. It separates an item's \emph{validity} from its \emph{usefulness}, selects under a hard token budget with a monotone submodular objective and a per-query optimality certificate, estimates validity from the store's own version history, and learns a linear policy from execution feedback. We prove a $\frac12(1-\frac1e)$ selection guarantee, a dynamic regret bound in which the price of expiring items is exactly the accumulated validity error, a linear lower bound for validity-blind learners, a rate--distortion bound for reversible consolidation, and a stationary store-size result. On retrieval-level metrics, with no LLM in the evaluation, \kanerva raises recall at 1024 tokens from \R{lme}{hybrid}{1024} to \R{lme}{kanerva}{1024} on LongMemEval-S and from \R{locomo}{hybrid}{1024} to \R{locomo}{kanerva}{1024} on LoCoMo against hybrid retrieval. On RepoMem, a new benchmark built from real repository histories, it cuts stale-fact errors on changed facts from \V{stalechanged-hybrid} to \V{stalechanged-kanerva} with no structured keys, at a cost of up to two points of recall. A cross-encoder reranker beats it at small budgets; feeding the reranker's signals into the objective, a post-hoc extension, closes that gap. Hybrid retrieval still wins at large budgets on LoCoMo, forgetting and aggressive consolidation cost recall, and online learning is fragile under reward noise. We report all of these cases in full.
\end{abstract}
